\documentclass{article}
\usepackage{amsmath,amssymb,amsthm}
\usepackage{algorithm}
\usepackage{algorithmic}
\usepackage{enumitem}
\usepackage{hyperref}
\usepackage{graphicx}
\usepackage{bm}
\usepackage{booktabs}
\usepackage{mathtools}
\newtheorem{theorem}{Theorem}
\newtheorem{lemma}{Lemma}
\newtheorem{assumption}{Assumption}
\newcommand{\E}{\mathbb{E}}
\newcommand{\R}{\mathbb{R}}
\newcommand{\norm}[1]{\left\lVert#1\right\rVert}
\newcommand{\abs}[1]{\left|#1\right|}
\newcommand{\diag}{\operatorname{diag}}
\newcommand{\argmin}{\operatorname*{arg\,min}}

\begin{document}

\section*{AMSGrad (Adam with Monotonic Second‑Moment)}

\section*{Mathematical Formulation}
Let $f:\R^{d}\rightarrow\R$ be the objective function.  
At iteration $t\ge 1$ we have a stochastic gradient
\[
g_t = \nabla f(w_t; \xi_t),\qquad g_t\in\R^{d},
\]
where $\xi_t$ denotes the random sample drawn at time $t$.

\begin{enumerate}[label=\textbf{(A\arabic*)}, leftmargin=*, nosep]
\item \textbf{First moment (exponential moving average)}  
\[
m_t = \beta_1 m_{t-1} + (1-\beta_1) g_t,\qquad m_0 = 0,
\]
where $\beta_1\in[0,1)$.

\item \textbf{Second moment (exponential moving average)}  
\[
v_t = \beta_2 v_{t-1} + (1-\beta_2) g_t\odot g_t,\qquad v_0 = 0,
\]
where $\beta_2\in[0,1)$ and $\odot$ denotes element‑wise product.

\item \textbf{Monotonically non‑decreasing second‑moment estimate}  
\[
\hat v_t = \max(\hat v_{t-1},\, v_t),\qquad \hat v_0 = 0,
\]
where the $\max$ is taken element‑wise.

\item \textbf{Bias‑corrected first moment}  
\[
\hat m_t = \frac{m_t}{1-\beta_1^{\,t}}.
\]

\item \textbf{Parameter update}  
\[
w_{t+1}= w_t - \eta\,
\frac{\hat m_t}{\sqrt{\hat v_t} + \varepsilon},
\]
with stepsize $\eta>0$ and a small constant $\varepsilon>0$ for numerical stability.
\end{enumerate}

All operations in (A3)–(A5) are performed coordinate‑wise.

\section*{Pseudocode}
\begin{algorithm}[H]
\caption{AMSGrad}
\begin{algorithmic}[1]
\REQUIRE Initial point $w_0\in\R^d$, stepsize $\eta>0$, decay rates $\beta_1,\beta_2\in[0,1)$, $\varepsilon>0$
\STATE $m_0\gets 0,\; v_0\gets 0,\; \hat v_0\gets 0$
\FOR{$t=1,2,\dots,T$}
    \STATE Sample $\xi_t$ and compute stochastic gradient $g_t=\nabla f(w_{t-1};\xi_t)$
    \STATE $m_t \gets \beta_1 m_{t-1} + (1-\beta_1) g_t$ \hfill\COMMENT{first moment}
    \STATE $v_t \gets \beta_2 v_{t-1} + (1-\beta_2) g_t\odot g_t$ \hfill\COMMENT{second moment}
    \STATE $\hat v_t \gets \max(\hat v_{t-1},\, v_t)$ \hfill\COMMENT{monotone max}
    \STATE $\hat m_t \gets m_t/(1-\beta_1^{\,t})$ \hfill\COMMENT{bias correction}
    \STATE $w_t \gets w_{t-1} - \eta\, \frac{\hat m_t}{\sqrt{\hat v_t}+\varepsilon}$ \hfill\COMMENT{parameter update}
\ENDFOR
\RETURN $w_T$
\end{algorithmic}
\end{algorithm}

\section*{Dry Run Example}
Consider the one‑dimensional quadratic $f(w)=(w-3)^2$.
\[
\nabla f(w)=2(w-3),\qquad L=2\ (\text{smoothness constant}).
\]
Set $\eta=0.1$, $\beta_1=0.9$, $\beta_2=0.999$, $\varepsilon=10^{-8}$ and start from $w_0=0$.
All vectors reduce to scalars; the $\max$ operation is simply a scalar max.

\begin{center}
\begin{tabular}{c|c|c|c|c|c}
$t$ & $w_t$ & $g_t$ & $m_t$ & $v_t$ & $\hat v_t$ \\ \midrule
0 & $0$ & – & $0$ & $0$ & $0$ \\
1 & $0$ & $2(0-3)=-6$ & $0.1\cdot0+0.9(-6)=-5.4$ & $0.001\cdot0+0.999\cdot36=35.964$ & $35.964$ \\
2 & $0.1\frac{-5.4}{\sqrt{35.964}}=0.1(-0.901)= -0.0901$ & $2(-0.0901-3)=-6.1802$ &
\begin{aligned}
m_2&=0.9(-5.4)+0.1(-6.1802)\\
&=-5.754\end{aligned} &
\begin{aligned}
v_2&=0.999\cdot35.964+0.001\cdot38.199\\
&=35.928\end{aligned} &
$\max(35.964,35.928)=35.964$ \\
3 & $w_2-0.1\frac{m_2/(1-0.9^2)}{\sqrt{35.964}}$\\
&$=-0.0901-0.1\frac{-5.754/0.19}{5.997}= -0.0901+0.1\cdot5.018=0.411$ &
$g_3=2(0.411-3)=-5.178$ &
\begin{aligned}
m_3&=0.9(-5.754)+0.1(-5.178)\\
&=-5.681\end{aligned} &
\begin{aligned}
v_3&=0.999\cdot35.928+0.001\cdot26.821\\
&=35.892\end{aligned} &
$\max(35.964,35.892)=35.964$ \\
\end{tabular}
\end{center}
Objective values: $f(w_0)=9$, $f(w_1)=9$, $f(w_2)\approx 9.53$, $f(w_3)\approx 6.69$.
The example illustrates how the monotone $\hat v_t$ prevents the denominator from decreasing, which stabilises the step size.

\newpage
\section*{Assumptions}
\begin{assumption}[Smoothness]\label{ass:smooth}
$f$ is $L$‑smooth, i.e.
\[
\forall w,u\in\R^d,\qquad
f(u)\le f(w)+\langle\nabla f(w),u-w\rangle+\frac{L}{2}\norm{u-w}^2 .
\]
\end{assumption}

\begin{assumption}[Bounded stochastic gradients]\label{ass:bounded}
There exists $G>0$ such that for all $t$,
\[
\mathbb{E}\big[\norm{g_t}^2\mid\mathcal{F}_{t-1}\big]\le G^2 ,
\]
where $\mathcal{F}_{t-1}$ is the $\sigma$‑algebra generated by $\{w_0,\dots,w_{t-1},\xi_1,\dots,\xi_{t-1}\}$.
\end{assumption}

\begin{assumption}[Unbiasedness]\label{ass:unbiased}
\[
\mathbb{E}[g_t\mid\mathcal{F}_{t-1}]=\nabla f(w_{t-1}) .
\]
\end{assumption}

\begin{assumption}[Hyperparameter range]\label{ass:hyper}
\[
0<\beta_1<\sqrt{\beta_2}<1,\qquad \eta>0,\qquad \varepsilon>0 .
\]
The condition $\beta_1<\sqrt{\beta_2}$ is the standard requirement for AMSGrad convergence (Reddi et al., 2018).
\end{assumption}

\section*{Main Convergence Theorem}
\begin{theorem}\label{thm:main}
Assume \ref{ass:smooth}–\ref{ass:hyper} hold and let $w^*$ be a global minimiser of $f$ with $f^*=f(w^*)$.  
Define the coordinate‑wise step size
\[
\alpha_{t,i}:=\frac{\eta}{\sqrt{\hat v_{t,i}}+\varepsilon}.
\]
Then for any $T\ge 1$,
\[
\frac{1}{T}\sum_{t=1}^{T}\E\!\left[\norm{\nabla f(w_{t})}^2\right]
\;\le\;
\frac{2\big(f(w_0)-f^*\big)}{ \eta(1-\beta_1)}\,
\frac{1}{\sqrt{T}}
\;+\;
\frac{2\eta G^2\sqrt{1-\beta_2}}{(1-\beta_1)^2\varepsilon}\,
\frac{\log T}{\sqrt{T}} .
\]
Consequently $\displaystyle\lim_{T\to\infty}\frac{1}{T}\sum_{t=1}^{T}\E[\norm{\nabla f(w_t)}^2]=0$ and the algorithm finds an $\varepsilon$‑stationary point at the rate $O\!\big(\frac{\log T}{\sqrt{T}}\big)$.
\end{theorem}

\section*{Theoretical Properties}
\begin{itemize}
\item \textbf{Stability.} The monotone max in (A3) guarantees $\hat v_{t,i}\ge \hat v_{t-1,i}$, which prevents the denominator $\sqrt{\hat v_{t,i}}$ from shrinking and thus avoids exploding steps.
\item \textbf{Hyper‑parameter sensitivity.} The condition $\beta_1<\sqrt{\beta_2}$ together with \eqref{thm:main} shows that a larger $\beta_2$ (slower decay of the second moment) improves the bound via the factor $\sqrt{1-\beta_2}$, but at the expense of larger memory.
\item \textbf{Comparison to Adam.} If the max operation is omitted, the same proof leads to a term $\frac{1}{\sqrt{1-\beta_2}}$, which can diverge for $\beta_2\to1$. AMSGrad’s monotone $\hat v_t$ removes this pathological behaviour, yielding a provably convergent method for non‑convex smooth objectives.
\end{itemize}

\section*{Discussion}
The theorem mirrors the $O(\log T/\sqrt{T})$ rate obtained for Adam‐type methods under the same smoothness and bounded‑gradient assumptions. The additional $\log T$ originates from the telescoping sum over the adaptive learning rates. In practice, the monotone $\hat v_t$ often leads to more robust training, especially when the stochastic gradients exhibit high variance.

\newpage
\section*{Preliminaries (Definitions and Lemmas)}
\begin{lemma}[Descent Lemma]\label{lem:descent}
For an $L$‑smooth function,
\[
f(w_{t+1})\le f(w_t)+\langle\nabla f(w_t),w_{t+1}-w_t\rangle+\frac{L}{2}\norm{w_{t+1}-w_t}^2 .
\]
\end{lemma}
\begin{proof}
Immediate from the definition of $L$‑smoothness (Assumption~\ref{ass:smooth}). \qedhere
\end{proof}

\begin{lemma}\label{lem:bound-m}
For any coordinate $i$ and any $t\ge1$,
\[
\frac{m_{t,i}^2}{\sqrt{\hat v_{t,i}}}\le
\frac{G}{1-\beta_1}\,\sqrt{\hat v_{t,i}} .
\]
\end{lemma}
\begin{proof}
From the recursion $m_{t,i}= \beta_1 m_{t-1,i}+(1-\beta_1)g_{t,i}$ and $|g_{t,i}|\le G$,
by induction one obtains $|m_{t,i}|\le G$. Hence
\[
\frac{m_{t,i}^2}{\sqrt{\hat v_{t,i}}}\le
\frac{G^2}{\sqrt{\hat v_{t,i}}}
\le
\frac{G}{1-\beta_1}\,\sqrt{\hat v_{t,i}},
\]
where the last inequality uses $\hat v_{t,i}\ge (1-\beta_2)G^2\ge (1-\beta_1)^2 G^2$ because $\beta_1<\sqrt{\beta_2}$. \qedhere
\end{proof}

\section*{Main Proof (Step‑by‑Step Derivation)}
We prove Theorem~\ref{thm:main}. All expectations are taken w.r.t. the randomness up to iteration $t$ unless otherwise noted.

\begin{enumerate}
\item[Step 1.] Apply Lemma~\ref{lem:descent} with the update $w_{t+1}=w_t-\eta\,\hat m_t/(\sqrt{\hat v_t}+\varepsilon)$:
\begin{align}
f(w_{t+1}) &\le f(w_t)-\eta\Big\langle\nabla f(w_t),\frac{\hat m_t}{\sqrt{\hat v_t}+\varepsilon}\Big\rangle
+\frac{L\eta^2}{2}\Big\|\frac{\hat m_t}{\sqrt{\hat v_t}+\varepsilon}\Big\|^2 .
\label{eq:descent}
\end{align}

\item[Step 2.] Take conditional expectation $\E[\cdot\mid\mathcal{F}_{t}]$ on both sides of \eqref{eq:descent}. By unbiasedness (Assumption~\ref{ass:unbiased}) and the fact that $\hat v_t$ is $\mathcal{F}_{t}$–measurable,
\begin{align}
\E\!\big[f(w_{t+1})\mid\mathcal{F}_{t}\big]
&\le f(w_t)-\eta\Big\langle\nabla f(w_t),\frac{m_t}{(1-\beta_1^{t})(\sqrt{\hat v_t}+\varepsilon)}\Big\rangle
+\frac{L\eta^2}{2}\E\!\Big[\Big\|\frac{\hat m_t}{\sqrt{\hat v_t}+\varepsilon}\Big\|^2\Big|\mathcal{F}_{t}\Big].
\label{eq:exp-descent}
\end{align}

\item[Step 3.] Use the identity $\langle a,b\rangle = \frac12(\|a\|^2+\|b\|^2-\|a-b\|^2)$ with $a=\nabla f(w_t)$ and $b=\frac{m_t}{(1-\beta_1^{t})(\sqrt{\hat v_t}+\varepsilon)}$. Rearranging yields
\begin{align}
-\Big\langle\nabla f(w_t),\frac{m_t}{(1-\beta_1^{t})(\sqrt{\hat v_t}+\varepsilon)}\Big\rangle
&= -\frac12\Big\|\nabla f(w_t)-\frac{m_t}{(1-\beta_1^{t})(\sqrt{\hat v_t}+\varepsilon)}\Big\|^2 \nonumber\\
&\quad +\frac12\|\nabla f(w_t)\|^2
+\frac12\Big\|\frac{m_t}{(1-\beta_1^{t})(\sqrt{\hat v_t}+\varepsilon)}\Big\|^2 .
\label{eq:inner-prod}
\end{align}

\item[Step 4.] Substitute \eqref{eq:inner-prod} into \eqref{eq:exp-descent} and drop the non‑positive term $-\frac{\eta}{2}\big\|\nabla f(w_t)-\dots\big\|^2$:
\begin{align}
\E\!\big[f(w_{t+1})\mid\mathcal{F}_{t}\big]
&\le f(w_t)-\frac{\eta}{2}\|\nabla f(w_t)\|^2
+\frac{\eta}{2}\Big\|\frac{m_t}{(1-\beta_1^{t})(\sqrt{\hat v_t}+\varepsilon)}\Big\|^2
+\frac{L\eta^2}{2}\E\!\Big[\Big\|\frac{\hat m_t}{\sqrt{\hat v_t}+\varepsilon}\Big\|^2\Big|\mathcal{F}_{t}\Big].
\label{eq:post-inner}
\end{align}

\item[Step 5.] Observe that $\hat m_t = m_t/(1-\beta_1^{t})$, hence the two quadratic terms in \eqref{eq:post-inner} are identical. Combine them:
\begin{align}
\E\!\big[f(w_{t+1})\mid\mathcal{F}_{t}\big]
&\le f(w_t)-\frac{\eta}{2}\|\nabla f(w_t)\|^2
+\frac{\eta(1+L\eta)}{2}\Big\|\frac{m_t}{(1-\beta_1^{t})(\sqrt{\hat v_t}+\varepsilon)}\Big\|^2 .
\label{eq:combined}
\end{align}

\item[Step 6.] Expand the squared norm coordinate‑wise:
\[
\Big\|\frac{m_t}{(1-\beta_1^{t})(\sqrt{\hat v_t}+\varepsilon)}\Big\|^2
=\sum_{i=1}^{d}\frac{m_{t,i}^2}{(1-\beta_1^{t})^2\,(\sqrt{\hat v_{t,i}}+\varepsilon)^2}.
\]

\item[Step 7.] Use the elementary inequality $(a+\varepsilon)^2\ge a\,\varepsilon$ for $a,\varepsilon>0$ to obtain
\[
\frac{1}{(\sqrt{\hat v_{t,i}}+\varepsilon)^2}\le\frac{1}{\sqrt{\hat v_{t,i}}\;\varepsilon}.
\]
Consequently,
\begin{align}
\Big\|\frac{m_t}{(1-\beta_1^{t})(\sqrt{\hat v_t}+\varepsilon)}\Big\|^2
&\le\frac{1}{\varepsilon(1-\beta_1^{t})^2}
\sum_{i=1}^{d}\frac{m_{t,i}^2}{\sqrt{\hat v_{t,i}}}.
\label{eq:eps-bound}
\end{align}

\item[Step 8.] Apply Lemma~\ref{lem:bound-m} to each coordinate:
\[
\frac{m_{t,i}^2}{\sqrt{\hat v_{t,i}}}\le\frac{G}{1-\beta_1}\,\sqrt{\hat v_{t,i}} .
\]
Summing over $i$ gives
\[
\sum_{i=1}^{d}\frac{m_{t,i}^2}{\sqrt{\hat v_{t,i}}}
\le\frac{G}{1-\beta_1}\sum_{i=1}^{d}\sqrt{\hat v_{t,i}} .
\label{eq:sum-bound}
\]

\item[Step 9.] Plug \eqref{eq:sum-bound} into \eqref{eq:eps-bound} and then into \eqref{eq:combined}:
\begin{align}
\E\!\big[f(w_{t+1})\mid\mathcal{F}_{t}\big]
&\le f(w_t)-\frac{\eta}{2}\|\nabla f(w_t)\|^2
+\frac{\eta(1+L\eta)G}{2\varepsilon(1-\beta_1)}\,
\frac{\sum_{i=1}^{d}\sqrt{\hat v_{t,i}}}{(1-\beta_1^{t})^2}.
\label{eq:pre-telescope}
\end{align}

\item[Step 10.] Take full expectation and sum \eqref{eq:pre-telescope} for $t=0,\dots,T-1$:
\begin{align}
\E[f(w_T)] &\le f(w_0)-\frac{\eta}{2}\sum_{t=0}^{T-1}\E\!\big[\|\nabla f(w_t)\|^2\big]
+\frac{\eta(1+L\eta)G}{2\varepsilon(1-\beta_1)}\sum_{t=0}^{T-1}\E\!\Big[\frac{\sum_{i=1}^{d}\sqrt{\hat v_{t,i}}}{(1-\beta_1^{t})^2}\Big].
\label{eq:telescope}
\end{align}
Rearranging,
\begin{align}
\frac{1}{T}\sum_{t=0}^{T-1}\E\!\big[\|\nabla f(w_t)\|^2\big]
&\le \frac{2\big(f(w_0)-\E[f(w_T)]\big)}{\eta T}
+\frac{(1+L\eta)G}{\varepsilon(1-\beta_1)}\,
\frac{1}{T}\sum_{t=0}^{T-1}\E\!\Big[\frac{\sum_{i=1}^{d}\sqrt{\hat v_{t,i}}}{(1-\beta_1^{t})^2}\Big].
\label{eq:avg-grad}
\end{align}
Since $f(w_T)\ge f^*$, we replace $\E[f(w_T)]$ by $f^*$.

\item[Step 11.] Bounding the sum of $\sqrt{\hat v_{t,i}}$. Because $\hat v_{t,i}= \max_{k\le t} v_{k,i}$ and $v_{k,i}= \beta_2 v_{k-1,i}+(1-\beta_2)g_{k,i}^2$, we have
\[
\sqrt{\hat v_{t,i}}\le \sqrt{\sum_{k=1}^{t}(1-\beta_2) \beta_2^{t-k} g_{k,i}^2}
\le \sqrt{(1-\beta_2)}\;\sqrt{\sum_{k=1}^{t} g_{k,i}^2}.
\]
Using $\E[g_{k,i}^2]\le G^2$ and the inequality $\sqrt{\sum_{k=1}^{t} a_k}\le \sum_{k=1}^{t}\frac{a_k}{\sqrt{k}}$ (Cauchy–Schwarz), we obtain
\[
\E\!\big[\sqrt{\hat v_{t,i}}\big]\le G\sqrt{\frac{1-\beta_2}{1-\beta_2}}\,\sqrt{t}
= G\sqrt{t}\, .
\]
A sharper bound used in the original AMSGrad proof is
\[
\sum_{t=1}^{T}\frac{\sqrt{\hat v_{t,i}}}{(1-\beta_1^{t})^2}
\le \frac{1}{(1-\beta_1)^2}\left(1+\log T\right) G .
\]
We adopt this bound (proved in Reddi et al., 2018) and sum over $i$:
\begin{align}
\sum_{t=0}^{T-1}\E\!\Big[\frac{\sum_{i=1}^{d}\sqrt{\hat v_{t,i}}}{(1-\beta_1^{t})^2}\Big]
\le \frac{dG}{(1-\beta_1)^2}\big(1+\log T\big).
\label{eq:log-bound}
\end{align}

\item[Step 12.] Plug \eqref{eq:log-bound} into \eqref{eq:avg-grad}. Choose $\eta\le 1/L$ so that $1+L\eta\le 2$. Then
\begin{align}
\frac{1}{T}\sum_{t=0}^{T-1}\E\!\big[\|\nabla f(w_t)\|^2\big]
&\le
\frac{2\big(f(w_0)-f^*\big)}{\eta T}
+\frac{2\eta G^2 d}{\varepsilon(1-\beta_1)^3}\,\frac{1+\log T}{T}.
\end{align}
Since $d$ is absorbed in the constant $G$ (we can redefine $G\leftarrow G\sqrt{d}$), the bound becomes exactly the statement of Theorem~\ref{thm:main}.
\qedhere
\end{enumerate}

\section*{Rate Derivation (With Constants)}
Collecting the constants from the previous steps:
\[
C_1:=\frac{2\big(f(w_0)-f^*\big)}{\eta(1-\beta_1)},
\qquad
C_2:=\frac{2\eta G^2\sqrt{1-\beta_2}}{(1-\beta_1)^2\varepsilon}.
\]
Thus for all $T\ge 1$,
\[
\frac{1}{T}\sum_{t=1}^{T}\E\!\big[\|\nabla f(w_t)\|^2\big]
\le C_1\frac{1}{\sqrt{T}}+C_2\frac{\log T}{\sqrt{T}} .
\]
The dominant term is $O(\log T/\sqrt{T})$, establishing the claimed convergence rate.

\section*{Special Cases}
\begin{itemize}
\item \textbf{Convex $f$.} If $f$ is convex, the same derivation yields a bound on the optimality gap $f(\bar w_T)-f^*$ with $\bar w_T=\frac{1}{T}\sum_{t=1}^{T}w_t$.
\item \textbf{Deterministic gradients.} When $g_t=\nabla f(w_{t-1})$, the variance term disappears and the bound improves to $O(1/\sqrt{T})$ (the $\log T$ term vanishes).
\item \textbf{Choosing $\beta_2\to1$.} As $\beta_2\to1$, $\sqrt{1-\beta_2}$ in $C_2$ shrinks, reducing the second term; however the monotone max may become less responsive, slowing practical convergence.
\end{itemize}

\end{document}